
\documentclass[11pt]{report}           %% ceci est un commentaire (apres le caractere %)
%\usepackage[ landscape, margin=1in]{geometry}

\usepackage[utf8]{inputenc}
\usepackage[english]{babel}
\setlength{\parskip}{1em}
\renewcommand{\baselinestretch}{1.6}

\usepackage{geometry}
 \geometry{
 a4paper,
 total={170mm,245mm},
 left=20mm,
 top=30mm,
 }
 
\usepackage{relsize}
 
\usepackage{amsthm}
\usepackage{array}
\usepackage{breqn}
\usepackage{amsmath}
\usepackage[ruled,vlined,linesnumbered,noresetcount]{algorithm2e}
\usepackage{algpseudocode}
\usepackage{dsfont}
\usepackage{tabu}
\usepackage{graphicx}
\graphicspath{ {img/} }
%\graphicspath{ {figures/} }
\usepackage{float}
\usepackage[justification=centering]{subfig}
\usepackage[colorlinks]{hyperref}
\usepackage{hyperref} 
\hypersetup{
    colorlinks=true,
    linkcolor=blue,
    filecolor=magenta,      
    urlcolor=cyan,
}
%\usepackage{algorithm}
%\usepackage{algpseudocode}
\usepackage{algorithm2e}
\usepackage{wrapfig}
\usepackage{listings}
\usepackage{caption}
%\usepackage{array,multirow,makecell}
%\setcellgapes{1pt}
%\makegapedcells
 \usepackage{multirow}
\usepackage{tabularx}
\usepackage{bm}
 %\renewcommand{\familydefault}{\sfdefault}
 
% \renewcommand{\familydefault}{\rmdefault}

\usepackage{lmodern}
\begin{document}



%\maketitle\thispagestyle{empty}

\begin{abstract}
Our work revolves around the recognition of both hands  gesture in Real time using the kinect sensor and machine learning. At first we acquired depth image of the kinect, based on the depth data we used its advantage to detect And perform gesture by  Considering that the minimum depth from the sensor  is the hand region, since the hand gestures are  usually held in front of  the body, this  allows us to track the hand and also it gives  the user more  mobility around the sensor rather than working in a certain closed and predefined range. After that we have developed an algorithm for gesture classification using SIFT (Scale Invariant Features Transform ) and SURF (Speeded up robust features ) descriptors  with Bag of Visual words approach to Quantify the keypoints Features, this approach consumes a great amount of time and memory affecting the efficiency and speed of our Robust HGR  system. however the second algorithm is Fourier descriptor which was developed to meet real time requirements.

These two approaches  were used to classify 16 Gestures using  three Machine Learning algorithms (Support vector Machine with K-means  and Nearest Neighbors classifier )    

\end{abstract}


%content
\tableofcontents
\listoffigures
\listoftables

\include{chapters/chapter0}


\chapter{LITERATURE SURVEY} \label{chapter1}

\section{introduction:} 

According to the World Federation of Deaf (WFD) \cite{wfd} there are 70 million deaf people in the world, whereas the world population in 2012 exceeded 7 billion, this number has got the power to make research in the field of machine vision  over  the  past  decade become rich with new scientific approaches.
A  more  recent   technology  for gesture   recognition  is to  incorporate   the  information   of object   distances,   or depths,   normally   captured   by a 3D camera,   or a set  of cameras  that  produce 3D images and that it is Technology that will be used in this project. It is also a contact-less  user interface,   in contrast   to glove based   devices  or handhold  remote  sensor.
\newline
With the kinect depth Sensor stream data, it is now possible to recognize and  translate sign language into speech and/or text in order to facilitate communication between normal hearing and hearing impaired people. This is usually called Sign Language Recognition (SLR) which  aim  to automatically understand the meaning of a sequence of signs from sign language without involvement of human interpreter \cite{27}. This chapter covers key aspects of SLR Literature.

\section{ Sign  Language Recognition}  

the SLR system can be divided into 4 important  steps, data acquisition and feature extraction and finally feature combination and sign recognition, following below, we will cover two important approaches of the first step ( data acquisition ) which are  :

\textbf{Device based }:
in  1991 Tomoichi Takahashi   developed a system for real time SLR  using Device based approach, this system required the user to wear gloves connected to the omputer for transmission of hand signal. Experiments showed that 34 of 46 gestures of the japanese Kana manual alphabet were recognized correctly \cite{29}
.
A wireless glove that was designed by Ryan Petters  \cite{31}in 2002, sensed hand movements that are involved in sign language and then transmitted them wirelessly to portable device, that displayed translated signs as lines of text.

\textbf{vision based }:
This approach  \cite{32} here  feature extraction was done using color camera for recognition of ASL. Two experiments were run: the first required from the signer to wear colored gloves,while the second used hand’s skin color for feature extraction\cite{32}. The extracted features were used as input for HMM algorithm that was used during sign recognition phase. The first experiment attained 99\% word accuracy while the second 92\%.
 
Two experiments were carried out also in \cite{33} on a SLR system that combined computer vision and HMM for recognition of sentence level continuous ASL. Both systems were tested with 40 words from ASL, where in the first system a single camera was mounted on desk to track user’s hands and achieved 92\% word accuracy, while in the second system the user had to wear a cap worm in which camera was mounted and the achieved accuracy was 97\%  \cite{33}.



The presentation of all the work related to SLR is not the intention of this section and more can be found on  \cite{27},\cite{36} and \cite{39}. 

The following section presents work that employed Kinect device for building SLR systems in which we are going to discover The main ideas behind our project 

\section{Gesture Recognition Using Kinect :}

in \cite{35} Zhou Ren et al presents a gesture recognition System using Kinect sensor accompanied with a novel distance metric called Finger-Earth Mover’s Distance (FEMD) to measure the hand shape dissimilarities, which achieves better performance than Shape Context, in terms of accuracy (a 93.9\% mean accuracy) and efficiency (0.5004s). 

A more recent approach followed by Daniel in \cite{47} uses Kinect skeletal tracking for features extraction. Other features were derived from skeletal tracking and finally 8-features were used for unique description of each sign. Sign recognition employs the nearest neighbor algorithm that tries to find the closest matching of performed sign within a group of stored sings  For a dictionary of 14 homemade signs the system achieves an accuracy of 95\% \cite{47}. 

in \cite{dardas} this paper presents an effective way of exploring depth information for hand gesture recognition. using  SIFT keypoint descriptor to extract local keypoint features over depth images,  which is an attempt for the first time in gesture recognition research. The recognition technique is inherently robust against scaling, rotation, and illumination conditions.  The features were  quantified  and clustered, and finally fed to  An SVM to be trained.\\


\textbf{ Insipired HGR system  of this project : }


\textbf{with \cite{dardas} and \cite{47} } we got inspired for our most work dedicated for this project, we take advantage of  depth and Skeletal tracking  data generated by Kinect to detect the hand and track it  in a range of [80 cm until 3 meters] long.
After detection of hand, we used SURF and SIFT descriptors for features detection and extraction then we  cluster these vectors via a bag-of-features (BOW) model, these features are then fed to an SVM with two different kernels linear and RBF kernel.

on the other hand unlike most systems like \cite{33} that use HMM,  we proposed and tested a method for static hand gesture recognition using  Fourier shape  descriptor which proved to be a very efficient way for feature extraction, it provided a set of invariant features for each image,  and provided a small number of features compared to local image descriptor, we found NN classifier to give faster and better results for Fourier descriptor.
\newpage
\section{Conclusion}

We can summarize The process used in SLR literature as the following:

Getting data from a sign performing can either be device based or/and vision approach, then features
involved in sign making must be extracted and finally combination of features must be analyzed 

\begin{figure}[H]
\centering
\includegraphics[width=0.8\textwidth]{img/SLR.PNG}
\caption{  Process of SLR }
\label{fig:SLR}
\end{figure}

In the next Chapter we will cover image descriptors as the second step of our HGR system,we will go through some of the important feature of these descriptor and how it can benefits our purposes in making Robust and accurate HGR System. 

\chapter{Descriptors}\label{descriptors}

\section{introduction}

Over the last decades, image feature detectors and descriptors have become popular tools in the computer vision community,  and they are being applied widely in different applications; image representation, image classification and retrieval, object recognition and matching  3D scene reconstruction, motion tracking   texture classification, robot localization, all rely on the presence of stable and representative features in the image. Thus, detecting and extracting the image features are vital steps for these applications.
Hence the need for a descriptor that can detect and extract features,
and since our HGR system needs a features' detector that can be stable in most cases which will fullfill our need to recognize the gesture performed by the user.


\section{Definitions and principles:}

\subsection{Global and local features}
the following definition was given by krig S in \cite{krig}. 

In the overall feature representation, the image is represented by one multidimensional feature vector which describes the image, in other terms, a single vector that is produced by the global representation method has different values that measure various aspects such as color, shape or texture. Practically we can extract multiple feature vectors from each image that allow us to compare two images. e.g if we want to differentiate sea pictures (blue) and forest pictures (green), a global color descriptor will provide us with different feature vectors of each category. In this context,  global features can be interpreted
as a particular property of image involving all pixels.
This property can be either texture, edges, colors histograms or even a specific descriptor extracted from some filters applied to the image \cite{h}.
On the other side, there is local features that distinctively represent some of the significant regions in the image, these regions remain invariant to viewpoint and intensity changes. 
This property can be color histograms, texture, edges or even a specific descriptor extracted from some filters
applied to the image. Thus, the image is represented
based on its local structures by a set of local feature descriptors which are extracted
from a set of image regions referred to as interest regions (i.e., keypoints), as illustrated in
Fig.\ref{fig:Ft1}

\begin{figure}[H]
\centering
\includegraphics[width=0.9\textwidth]{img/features.PNG}
\caption{ Global and local image features representation }
\label{fig:Ft1}
\end{figure}

In general, the kind of features we choose to use is greatly dependable on the applications on hand. Developers tend to use the most discriminative ones. e.g a person with big nose compared to smaller eyes, and a person with small nose compared to bigger eyes can have similar mug shot in terms of histogram or intensity distribution.\\


Then, local features or the global pattern distilled from local feature clusters seem to be more discriminative. Whereas, for very large datasets in the web scale image indexing application, it is appropriate to consider global features, as they are useful in applications where a rough segmentation of the object of interest is available.
The advantages of global features lay in the fact that they are much faster, are compact and are easy to compute and generally require small amounts of memory.\\ However, the global representation suffers from well known limitations, in particular they are not invariant to significant transformations and are sensitive to clutter and occlusion.\\In some applications, like copy detection, most of the illegal copies are very  similar to the original; they have only suffered from compression, scaling or limited
cropping. In contrast, the advantage of local features is their superior performance \cite{i}. Meanwhile, using local features for large scale image search have much higher performance than global features provide \cite{k}.\\ Besides, as the local structures are more distinctive and stable than other structures in smooth regions, it is probable to be more useful for image matching and object recognition.\\ However, they usually require a significant amount of memory because the image may contain hundreds of local features.\\ Researchers suggest aggregating local image descriptors into a very compact vector representation and optimizing the
dimensionality reduction of these vectors as a solution for this problem. \cite{k}


\subsection{Characteristics of feature detectors}

Tuytelaars and Mikolajczyk \cite{MM} define a local feature as \textit{ “it is an image pattern
which differs from its immediate neighborhood”.} Thus, they consider the purpose of
local invariant features is to provide a representation that allows to efficiently match
local structures between images. That is, we want to obtain a sparse set of local
measurements that capture the essence of the underlying input images and encode
their interesting structures. To meet this goal, the feature detectors and extractors
must have certain properties, keeping in mind that the importance of these properties
depends on the actual application settings and compromises need to be made. The
following properties are important for utilizing a feature detector in computer vision
applications:
\begin{itemize}
\item Robustness, the feature detection algorithm should be able to detect the same feature
locations independent of scaling, rotation, shifting, photometric deformations,
compression artifacts, and noise
\item Repeatability, the feature detection algorithm should be able to detect the same
features of the same scene or object repeatedly under variety of viewing conditions.
\item Accuracy, the feature detection algorithm should accurately localize the image
features (same pixel locations), especially for image matching tasks, where precise
correspondences are needed to estimate the epipolar geometry.
\item Generality, the feature detection algorithm should be able to detect features that
can be used in different applications.
\item Efficiency, the feature detection algorithm should be able to detect features in new
images quickly to support real time applications.
\item Quantity, the feature detection algorithm should be able to detect all or most of the
features in the image. Where, the density of detected features should reflect the
information content of the image for providing a compact image representation.
\end{itemize}

\textbf{in this project We are going to use Local features( descriptors) for its efficiency and scale, rotation  invariant }\\
\section{Spectra descriptors :}
\subsection{SIFT - Scale Invariant Feature Transforms} \label{siftSection}

The Scale Invariant Feature Transform (SIFT) developed by Lowe \cite{lowe}  it is basically the most powerful tool in vision field for detecting and description of interest points in an image, and to identify similar keypoint in different images ( matching ). The idea is these feature provide invariance to scale, rotation, affine distortion, noise,illumination.
the author \cite{lowe} also insisted on using several SIFT descriptors together gives additional invariance to occlusion and clutter. Here are the stages an image passed in order to get a sift descriptor :


\begin{enumerate}
\item Scale-space extrema detection :
The first stage of computation must search over all
scales and image locations, but it can be implemented efficiently by using a difference of-Gaussian
function to identify potential interest points that are invariant to scale and orientation.
\item Keypoint localization :
At each candidate location, a detailed model is fit to determine
location, scale, and contrast. Keypoints are selected based on measures of their
stability.
\item Orientation assignment :
One or more orientations are assigned to each keypoint
location based on local image properties. All future operations are performed relative to the assigned orientation, scale, and location for each feature, providing invariance
to these transformations.
\item Keypoint descriptor  :
The local image gradients are measured at the selected scale
in the region around each keypoint, and transformed into a representation that allows
for local shape distortion and change in illumination.
\end{enumerate}



\subsection{Scale-Space :}

Space scaling  is an important theory in computer vision.
It was developed by Tony Lindeberg in 1994, it allows to analyze an image with a multi-scale approach and to highlight different  structural  sizes in an image.

To create a scale space, SIFT takes the original image, and generates blurred images progressively by applying Gaussian blur.

\begin{align} 
 G(x, y,\sigma )  =\frac{1}{2\pi\sigma^2} e^{\frac{-( x^2 -  y^2)}{2\sigma^2}}
\end{align}


after  blurring images it resizes the original image at half its size and generates again blurred images and so on and so forth.
Images of the same octave are all blurred progressively  with a $\sigma$ (blurring coefficient). An example of the space construction of SIFT scales is presented in Figure \ref{fig:scale}.



\begin{figure}[H]
\centering
\includegraphics[scale = 0.25]{img/scale_space.pdf}
\caption{ constructing scale space.}
\label{fig:scale}
\end{figure}

\subsection{Detection of Extrema}
In the previous step, we constructed the scale space which consists of image smoothing on several scales and at different sizes. The idea now consists of making a deference of Gaussians (DoG), between two consecutive images of the same octave to obtain DoG images. These [DoGs]  are very useful for the detection of stable Features.


The difference between two consecutive images smoothed by a Gaussian filter constitutes a good approximation of the LoG (Laplacian of Gaussian), a high-pass filter is used to highlight  details that have a fast variation to the brightness  as well as  making the contours  of the objects  visible too, as shown in Figure \ref{fig:sift2}


\begin{figure}[H]
\centering
\includegraphics[width=0.5\textwidth]{img/scaleanddog.png}
\caption{ Detection of maxima and minima of the difference-of-Gaussian.}
\label{fig:sift2}
\end{figure}


The keypoints founded by DOG (that are stable) has  locals extremas through the deferent scales. Each pixel of the DoG images is then compared to its 26 neighbors, 8 neighbors in the same scale and 9 neighbors on the two neighboring scales that surround it.

\begin{figure}[H]
\centering
\includegraphics[width=0.5\textwidth]{img/sift2.PNG}
\caption{ Locate maxima/minima in DoG images }
\label{fig:sift4}
\end{figure}
If the pixel  is a local extremas, i.e, if it is greater or lesser  than its neighbors then it is selected as a candidate point for the next steps. This Comparison operation can only be performed for images that have neighboring images above and below.

\subsection{localization  of keypoints and elimination of irrelevant points:}

In general, the extremas detection step produces a large number of keypoint candidates, some of which are unstable;  in addition, their locations in particularly at  the higher scales (i.e. in the upper octaves of the pyramid where resolution is lower) remains approximate. In this step  we will precisely locate the keypoint (the detected extremas), and eliminate the points that have a low contrast and are located on the edges.

The idea is to calculate two gradients at the keypoint, both perpendicular to each other. Based on the image around the keypoint, three possibilities exist. The image around the keypoint can be:
\begin{itemize}

\item A flat region: If this is the case, both gradients will be small.
\item  An edge :  one gradient will be big (perpendicular to the edge) and the other will be small (along the edge)
\item A "corner":  both gradients will be big.
\end{itemize}
Corners are great keypoints. So we want just corners. If both gradients are big enough, we let it pass as a key point. Otherwise, it is rejected.




\subsection{Orientation Assignment }


After these steps we have already eliminated many Keypoints which were likely to be irrelevant to the implementation of the algorithm. Now, it’s a matter of attributing to each detected keypoint 
an orientation by collecting the gradient directions and magnitudes around this keypoint, then figure out the most prominent orientation (note it can be more than just one orientation, either way we assign the most prominent orientations to the detected keypoint). The next figure shows the magnitude and gradient orientations of a keypoint region.

\begin{figure}[H]
\centering
\includegraphics[width=18.5cm, height=13cm]{img/dog_hand.png}
\caption{ calculating Gradient and magnitude in the keypoints neighbor }
\label{fig:orientation}
\end{figure}



The magnitude and orientation is calculated for all pixels around the keypoint. Then, a histogram is created for this.
In this histogram, the 360 degrees of orientation are broken into 36 bins (each 10 degrees). if  the gradient direction at a certain point (in the "orientation collection region") is 18.759 degrees, then it will go into the 10-19 degree bin, and the amount that is added to the bin is proportional to the magnitude of gradient at that point.
then we add a component to the feature  vector of the keypoint that is  defined by (x, y, $\sigma$, $\theta$) where $\theta$ is the orientation of the highest peak in the histogram.

Also, any peaks above 80\% of the highest peak are converted into a new keypoint. This new keypoint has the same location and scale as the original. But it's orientation is equal to the other peak.

\begin{figure}[H]
\centering
\includegraphics[width=0.8\textwidth]{img/sift_histogram.jpg}
\caption{histogram of magnitude and orientation keypoint region}
\label{fig:histogram}
\end{figure}



In the figure above, you see the histogram peaks at 20-29 degrees. So, the keypoint is assigned orientation 3 (the third bin)

\subsection{keypoint descriptor:}
The last step is to generate a very characteristic [imprint] for each key point. This imprint, a [vector] aims to identify the keypoint, it is unique and tolerant respectfully to the comparison operation. In fact, things are never exactly the same when comparing two different images.
To do this, a 16x16 size window is created around the keypoint. This window is divided into windows of 4x4 size.

\begin{figure}[H]
\centering
\includegraphics[width=0.8\textwidth]{img/sift4.jpg}
\caption{feature generation}
\label{fig:feature}
\end{figure}


to the first bin. 45-89 add to the next bin. And so on. And the amount added to the bin depends on the magnitude of the gradient. 

Doing this for all 16 pixels within a (4 $\times$ 4) block,we  would gather  16 totally random orientations into 8 predetermined bins.  So you end up with 4$\times$ 4 $\times$8 = 128 numbers. These 128 numbers form the "feature vector". This keypoint is uniquely identified by this feature vector.



\subsection{ SURF ( Speeded Up Robust Features ) }\label{surf}

The detector and descriptor SURF is mainly known for its fast calculations. The comparative study shows the superiority of the SURF detector/descriptor in comparison with SIFT from a point of view of its performances in the execution time and its robustness to luminosity changes [Juan and Gwum, 2009]. In fact, SURF is inspired from SIFT algorithm which is the precursor in the extraction domain of Keypoints. However, the SIFT algorithm has the disadvantage of being slow, which blocks its application in real time.

The approach proposed by SURF uses an approximation of the Hessian matrix in order to detect blobs type structures.  The determinant  of this matrix  can be used to classify the maxima and minima of a function by the second order derivative,  the sign of the  $\det$ (H) can give an idea about local extremum which reprensents the keypoints.

$$ \det (H) =\frac{\partial^2 f \partial^2 f} {\partial x^2 {\partial y^2}}  \bm{-}  (\frac  {\partial^2 f}{\partial x.\partial y})^{2}$$


$
\left\{\begin{matrix}
\ \ \ \ \ \ \det(H) < 0 \Rightarrow  Eijenvalues\ have\ different\ signs\ and\ hence\ the\ point\ is\ not\ a\ local\ extremum.\\
\det(H) > 0 \Rightarrow  both\ Eijenvalues\ have\ same\ signs\ and\ hence\ the\ point\ is\ an\ extremum. 
\end{matrix}\right.
$


Here the hessian matrix is computed at a function of both space X=(x,y) and scale $\sigma$ :



\begin{gather}
 \mathcal{H}(X,\sigma) =
\begin{bmatrix}
                 {L_{xx}(X,\sigma)}  && {L_{xy}(X,\sigma)} \\
                 {L_{xy}(X,\sigma)} && {L_{yy}(X,\sigma)}
\end{bmatrix}
\end{gather}

Where ${L_{xx}(x,\sigma)}$ is the convolution of the Gaussian second order derivative with the image I in point x, and similarly for ${L_{xy}(x,\sigma)}$ and ${L_{yy}(x,\sigma)}$.
First convolution, then second order derivative, the approximation  of Laplacien of Gaussian  (LoG) using box filters for  the respective kernels :

\begin{figure}[H]
\centering
\includegraphics[width=0.55\textwidth]{img/surf4.PNG}
\caption{${L_{yy}(x,\sigma)}$ and${L_{xy}(x,\sigma)}$ Discretized
Gaussians and the approximations $D_{yy}$ and $D_{xy}$}
\label{fig:surf1}
\end{figure}


These approximate second order Gaussian derivatives and can be evaluated at a very low computational cost using integral images, and this is part of the reason why SURF is fast.

Now we can represent the determinant of the Hessian (approximated) as:

\begin{figure}[H]
\centering
\includegraphics[width= 0.4\textwidth]{img/det.jpeg}

\label{fig:surfdet}
\end{figure}

\subsection{Surf scale space }
Just like SIFT, SURF also created its own scale space.
The scale space is generally implemented as image pyramid (octaves). These images are successively smoothed with a Gaussian and then are sub-sampled in order to attend a high level of the pyramid, (from the largest to the smallest). In  the other hand, SURF proceeds differently in this, In fact, instead of applying successively the same process of Gaussian filtering of each image in  each octave (see SIFT), SURF uses the box-filters of different sizes directly on the integral image. Thus the scale space is built in enlarging the filter rather than iteratively reducing the image size. This enables first to reduce the calculation time, second it allows to avoid aliasing due to sub-sampling of the image.
\begin{figure}[H]
\centering
\includegraphics[width=0.6\textwidth]{img/surf2.png}
\caption {The left image represents the classical method with sub-sampling and same filtering. In the image on the right the filters are of various sizes.}
\label{fig:surf2}
\end{figure}


\subsection{Orientation assignment : }

We first calculate the Haar wavelet responses in x and y direction within a circular neighborhood of radius 6s around the interest point, with s the scale at which the interest point was detected. We calculate the sum of vertical and horizontal wavelet responses in a scanning aria, then change the scanning orientation (add $\pi$/3), and recalculate, until we find the orientation with largest sum value, this orientation is the main orientation of feature descriptor.   The size of the Haar filter is 2s. This process makes the descriptor invariant to rotation.

\begin{figure}[H]
\centering
\includegraphics[width=0.6\textwidth]{img/surf3.jpg}
\caption{descriptor computation }
\label{fig:surf3}
\end{figure}

\subsection {Descriptor component : }



The descriptor calculation is done on a squared neighborhood of 20s side, and oriented according to  the largest sum value. The repartition of this neighborhood is identical with the one performed for the SIFT descriptor, i.e 4x4 quadrants.  The dx and dy Haar wavelets are calculated inside each quadrant for a set of 5x5 points. Each quadrant is then characterized by a four components vector.

V=( $\sum dx , \sum \left | dx \right | ,\sum dy , \sum \left | dy \right |$)



\begin{figure}[H]
\centering
\includegraphics[width=0.5\textwidth]{img/surf5.png}
\caption{computation of dx and dy of Haar  }
\label{fig:surf5}
\end{figure}




The concatenation of these vectors lead to the SURF descriptor of 64 components,
our final descriptor is a 64-D vector. (In Sift, our descriptor is 128-D vector, so this is part of the reason that SURF is faster than Sift.)


\section{Differences and preferences :}
The two main advantages of SURF over SIFT is that SURF uses Laplacian of Gaussian so as to have a distinction
between background and foreground features, and secondly, SURF uses only 64 dimensional vector compared to 128 dimensional vector for SIFT. This helps in fast feature computation and also the quick matching capability.
The two different steps used by SURF to determine the local descriptor vectors are A. Keypoint detection B. Keypoint descriptor, which were  explained above \ref{surf} and \ref{siftSection}.

This makes SURF as our best candidate for gesture recognition  since it offers high number of stable features against scale invariance and rotation invariance, of course lesser than SIFT but it does not worth the computational operations of up scaling pyramid used in sift Nor the \textbf{valuable }time consumed.

Later on we will experimentally demonstrate that SURF achieved  better accuracy than SIFT in hand gesture recognition. However Even if SURF is less expensive than SIFT but still computationally expensive which is bad for our purposes,since what describe a good HGR is its ability to interpret  and interact with the signer performer in real time, That's where Fouriere descriptor appears to give accurate and a very fast feature extractor which will meet our requirement for a robust and real time application.


\section{Fourier shape descriptor } \label{FDT}
The Fourier transformed coefficients from the Fourier
shape descriptors represent the shape of the object in
the frequency domain. The general features of the shape can
be found in the lower frequency descriptors, while the higher
frequency descriptors contain information about the shape
details \cite{20}, before applying Fourier transform on the shape boundary, it is first normalized for matching purposes; this is done by sampling the boundary of each shape to have the same number of data points. The larger the number of sampled points the more details in the representation of the shape,  this results in more accurate matching. While a smaller number of sampled points reduce the accuracy of the matching results but on the other hand it will improve the computational efficiency.
After normalization we  apply Fourier transform to the
shape signature. A shape signature is any 1-D function
representing 2-D areas or boundaries. The shape signature we
used is complex coordinates. A complex coordinates function
is the complex form of the boundary coordinates \cite{20}.
\begin{gather}
    z_{i} = x_{i}+j y_{i} , i\in [1,N]
\end{gather}
For each shape we select N points with equal point
sampling. In order to facilitate the use of Fast Fourier
Transform (FFT), the number of sampled points is chosen to
be power of two. Assuming the number of sampled points is
N the Fourier transform gives N Fourier coefficients Cl. The
coefficients are usually called Fourier descriptors of the shape. 
\begin{gather}
C_{l}= \sum_{i=0}^{N-1} z_{i}e{\frac{-j 2\pi il}{N}} , l = 0,...,N-1
\end{gather}
The magnitude of the Fourier
transform of this set forms a unique shape signature, which can be used for generalized
gesture classification.
In addition, this descriptor is rotationally invariant. Shifts in the silhouette contour
points, which is the cause of rotation, will be appear as phase delays in frequency
domain. However, since only the magnitude of the Fourier coefficients is considered,
the phase (or equivalently, the rotation) is ignored. So, this method is rotationally
invariant while remaining computationally fast\\


\textbf{Generating the Fourier Descriptor: }
\begin{enumerate}
    \item $X_{c} = \frac{1}{N}\sum_{n=0}^{N-1} x(i) , Y_{c} = \frac{1}{N}\sum_{n=0}^{N-1} y(i)  = 1$ , where N is number of hand object pixels 
    \item Take the magnitude of the N point DFT of these points :\\
              \textbf{ abs(FT{r[n]}) = a[m], for m = 0..N-1 }.
    \item Normalize the Fourier coefficients by the DC value (Scale Invariance).
    \item Keep the first 7 normalized coefficients (skip DC, which is always 1).

\end{enumerate}


This is the Fourier descriptor for a single shape.\\

\textbf{Make a dictionary : } For a set of images of the same gesture, compute the average Fourier shape descriptor and add it to the dictionary with the same label.  Repeat for all desired gestures.

\textbf{Classification  : }Compute the Fourier descriptor for each new sample, compare it with each stored gesture in the dictionary using  Euclidean distance measure.  The label of the minimum distance is the desired gesture( nearest neighbor happens to be the most efficient method of classification )


 \section{ Conclusion : }
 In this chapter we went through two types of descriptors:
\textbf{the first} are  Local descriptors SIFT and SURF  descriptors that typically involve more intense computations and algorithms, often requiring floating point calculations, and may consume considerable memory, but generally they give a high  accuracy additionally to their stable performance even against scale variation and Rotation variation too, the following table is summarizing the fifferences of these two descriptors :

\begin{table}[H]
\centering

\begin{tabular}{l|l|l|ll}
\cline{2-3}
\textbf{}                                                                                     & \textbf{SIFT}                                                                                                                                                     & \textbf{SURF}                                                                                                                             &  &  \\ \cline{1-3}
\multicolumn{1}{|l|}{\textbf{\begin{tabular}[c]{@{}l@{}}Keypoint\\ Detection\end{tabular}}}   & \textbf{\begin{tabular}[c]{@{}l@{}}Different scale image\\ convoluted with\\ Gaussian function\end{tabular}}                                                      & \textbf{\begin{tabular}[c]{@{}l@{}}Original Image is\\ convoluted with\\ Different scale box\\ filter\end{tabular}}                       &  &  \\ \cline{1-3}
\multicolumn{1}{|l|}{\textbf{\begin{tabular}[c]{@{}l@{}}Keypoint\\ Description\end{tabular}}} & \textbf{\begin{tabular}[c]{@{}l@{}}Gradient amplitude of\\ a square area is\\ calculated with maximum gradient\\ strength as the main\\ direction\end{tabular}} & \textbf{\begin{tabular}[c]{@{}l@{}}A Haar Wavelet response\\  is used to calculate \\ each sector in a \\ circular area\end{tabular}} &  &  \\ \cline{1-3}
\multicolumn{1}{|l|}{\textbf{Dimensions}}                                                     & 128                                                                                                                                                               & 64                                                                                                                                        &  &  \\ \cline{1-3}
\end{tabular}
\caption{COMPARISON OF SIFT AND SURF}
\label{surfvssift}
\end{table}

\textbf{The second } are basis space descriptors, a basis space is composed of a set of functions, the basis functions, which are
composed together as a set, such as a series like the Fourier series. 

Fourier descriptors represent feature data as sine and cosine terms, which can be
observed in a Fourier power spectrum. The Fourier series, Fourier transform, and fast
Fourier transform are used for a wide range of signal analysis, including 1D, 2D, and 3D
problems. No discussion of image processing or computer vision is complete without
Fourier methods.

With this summary we are done from Feature detection / extraction and we will move on to the next phase of building  our HGR System, which is  feature classification using Machine learning algorithms that's going to be the goal of next chapter.

\chapter{Machine Learning} \label{ML}

\section{Introduction}

" We use AI so much that we have stopped thinking of it as AI, ”  says Nello Cristianini,  professor of artificial  intelligence at the university of Bristol.

Nowadays we make many decisions based on our past experiences,  or based on what we expect in the future, either way we sure take a lot of bad decisions and make a lot of mistakes, Machines ( computers ) Can be a powerful tool for this type of tasks with higher prediction ( about future ) and better accuracy, and that's when machine learning comes, which is actually a process  that teaches computers to tackle enormous amounts of data, and pull out patterns that can represent a model able to  make decisions, but these models are still yet in need of human supervision to evaluate it, and help it overcome the complexities hidden in data .

In this chapter we will first define supervised learning and presents one of the powerful application of such type of learning ( supervised ) SVM algorithm ( support vector machine )  and nearest neighbor algorithm, and then define the second type unsupervised learning   with a review K-means algorithm, these algorithm  proved  to be accurate and robust  for our HGR system as it's going to be shown in  chapter  [\ref{hgr}]

\section{Supervised machine learning}

The goal of supervised learning is to approximate the mapping function so well that when we have new input data (x) then we can predict the output variables (Y) for that data. When (Y) categorical,the problem is known as <<classification>>, and when (Y) represents floating values, in this case the process is knows as <<regression>> .

Some popular examples of supervised machine learning algorithms are:
Linear regression for regression problems.
Support vector machines for classification and regression problems.

\section{Support vector machines (SVM) }\label{sec:svm}
It's a supervised machine learning algorithm, the intuition  behind it  is to find a hyperplane that best divides a dataset into two classes, a good separation is achieved by the hyperplane, that has the largest distance to the nearest training-data point of any class (so-called functional margin), since in general the larger the margin the lower the  error of the classifier.

But first let us define what is a hyperplan;  it's a line that  linearly separates and classifies a set of data. As the fig \ref{122} shows, our hyperplan here is the line that separate the green data from red data .

\begin{figure}[H]
\centering
\includegraphics[width=0.3\textwidth]{img/svm1.png}
\caption{example of linear Hyperplan seperating Green data from Red data }
\label{122}
\end{figure}


The above is a classic example of a linear classifier, the hyperplan should be optimal for our classification task which turns to maximize the margin, which is basically the distance between the hyperplane,  and the nearest data point from either set,  the greater this margin is the better for the classifier to make  describe data .But data is not always linearly separable as in real world problems we usually deal with linearly non separable kind of data as shown  in the figure \ref{123 }, 

\begin{figure}[H]
\centering
\includegraphics[width=0.3\textwidth]{img/svm2.png}
\caption{Non linear Example of data  }
\label{123 }
\end{figure}
For this reason, it was proposed that the original finite-dimensional space be mapped into a much higher-dimensional space, presumably making the separation easier in that space, this is often called the kernel trick and it uses a set of mathematical functions (kernels). Using a higher dimension we can  linearly  seperate the data, thus, instead of constructing the complex curve (left schematic), all we have to do is to find an optimal line that can separate ( right schematic ) .



\begin{figure}[H]
\centering
\includegraphics[width=0.5\textwidth]{img/SVMIntro3.png}
\caption{Non linear separated using kernel (higher dimension) }
\label{124 }
\end{figure}


\subsection{Hard-Margin SVM}
The SVM technique is a classifier that finds a hyperplane or a function $g(x) = {\omega}^T +  b$   that correctly separates two classes with a maximum margin.The figure below shows a separating hyperplane corresponding to a hard-margin SVM (also called a linear SVM).

\begin{figure}[H]
\centering
\includegraphics[width=0.5\textwidth]{img/hardmargin.PNG}
\caption{ Hard-maximum-margin separating hyperplane. }
\label{125 }
\end{figure}


\subsection{Soft-Margin SVM}
SVM is classifying most of the data correctly, while allowing the model to misclassify a few points in the vicinity of the separating boundary

\begin{figure}[H]
\centering
\includegraphics[width=0.5\textwidth]{img/softmargin.PNG}
\caption{  A few misclassifications, as part of soft-margin SVM . }
\label{126 }
\end{figure}


\subsection{Kernels}

In most Real cases, the optimal boundary is not linear,  the idea of NON linear SVM is to map the data into a higher dimensional space where data can linearly separable, this is often called the kernel trick . A kernel should be a Hermitian and positive semidefinite matrix and needs to satisfy Mercer’s theorem, which translates into evaluating the kernel or Gram matrix on all pairs of data points as positive and semi-definite, forming:

$K(\small{x,u})=\sum\limits_{r} \phi(x)\phi(u) $ .

where $\phi(x) $ belongs to the Hilbert space.
\newline
$$\int \int K(x,u) g(x) g(u) du dx \geq 0 \ \ \ \ \forall g(x) , \ where \int g^2(x) dx < +\infty$$
\\
\newline 
Some popular kernel functions include: 


\begin{itemize}

  \item  Linear kernel  : 
  \begin{align*}
  \mathlarger {K(x,u)=  x^{T}.u  }
   \end{align*}
  \item Polynomial function: 
  \begin{align*}
  K(x,u)=(ax^{T}u + c)^{q} ,\, q>0.
  \end{align*}
  \item Gaussian radial basis function (RBF): 
\begin{align*} 
 K(x,u) = e^{\mathlarger{\frac{\mathbf{-\| x -  u\|^2}}{\mathbf{\sigma^2} } }}
\end{align*}
   
  \item  Hyperbolic tangent (sigmoid):
    \begin{equation}
      K(\mathbf x, \mathbf u) = \tanh(\mathbf {\beta} \mathbf x^{T} \mathbf u - \mathbf{\delta})^p
    \end{equation}
    
    \item Laplacian radial basis function : 
    
    \begin{align*} 
      K(x,u) = e^{\mathlarger{\frac{-\| x -  u\|}{\sigma} } }
     \end{align*}
    
\end{itemize}

Kernel selection is heavily dependent on the data specifics. For instance, the linear kernel the simplest
of all is useful in large sparse data vectors. However, it ranks behind the polynomial kernel, which avoids zeroing the Hessian. \\The polynomial kernel is widely used in image processing, the Gaussian and Laplace RBFs are general-purpose kernels that are mostly applied in the absence of prior knowledge. \\A kernel matrix that ends up being diagonal indicates that the feature space is redundant and that another kernel should be tried after feature reduction.\\\textbf{
Note} that when kernels are used to transform the feature vectors from input space to kernel space for linearly non-separable datasets, the kernel matrix computation requires massive memory and computational resources for big data . 



\section{Nearest neighbour based classifiers}

One of the simplest decision procedures that can be used for classification is the
nearest neighbour (NN) rule. It classifies a sample based on the category of its nearest
neighbour.When large samples are involved,it can be shown that this rule has a
probability of error which is less than twice the optimum error
—hence there is less
than twice the probability of error compared to any other decision rule. The nearest
neighbour based classifiers use some or all the patterns available in the training set
to classify a test pattern. These classifiers essentially involve finding the similarity
between the test pattern and every pattern in the training set.

\subsection{Nearest neighbour algorithm}
The nearest neighbour algorithm assigns to a test pattern the class label of its closest
neighbour

% Insert the algorith

\begin{algorithm}[H]
\SetAlgoLined

 Initialization\;
 $ A=(X_{1},\omega_{1}),(X_{2},\omega_{2}),....,(X_{n},\omega_{n}) \ // A\ is\ the\ set\ of\ N\ training\ pattern $\\
 
 
$ //where\  X_{i}\ is\ of\ dimension\ n\ and\  \omega_{i}\ is\ the\ class\ label\ of\ the\ i^{th}\ pattern $\\
 
 
 $d(u,x) = \sqrt{\sum_{k=0}^{N} (u_{i}-x_{i})^{2}}  // \ Euclidean\ distance\ measure $\\ 

 \While{ i \neq N }{ // N samples
 
  d(X_{t},X_{i}) = min {D(X_{t},X_{i})}\; 
 
 }
 
 $ X_{t} = \omega_{k} $ \\
 
 $ pattern\ X_{t}\ is\ assigned\ to\ a\ class\ \omega_{k}\ associated\ with\  X_{k} $
 
 \caption{Algorithm for NN}
\end{algorithm}



\subsection{K-nearest neighbour (kNN) algorithm}
 the k-nearest neighbors algorithm (k-NN) is a non-parametric method used for classification and regression.  In both cases, the input consists of the k closest training examples in the feature space. The output depends on whether k-NN is used for classification or regression:
\begin{itemize}
\item In k-NN classification, the output is a class membership. An object is classified by a majority vote of its neighbors, with the object being assigned to the class most common among its k nearest neighbors (k is a positive integer, typically small). 
\item In k-NN regression, the output is the property value for the object, this value is the average of the values of its k nearest neighbors.
\end{itemize}

K-NN is a type of instance based learning, or lazy learning, where the function is only approximated locally and all computation is deferred until classification. The k-NN algorithm is among the simplest of all machine learning algorithms.\\ Both for classification and regression, it can be useful to assign weight to the contributions of the neighbors, so that the nearest neighbors contribute more to the average than the more distant ones. For example, a common weighting scheme consists in giving each neighbor a weight of 1/d, where d is the distance to the neighbor.\\ The neighbors are taken from a set of objects for which the class (for k-NN classification) or the object property value (for k-NN regression) is known. This can be thought of as the training set for the algorithm, though no explicit training step is required.



%%%%%%%%%%%%%%%%%%%%%%%%%
\section{Unsupervised learning}

In supervised learning, the aim is to learn a mapping from the input to
an output whose correct values are provided by a supervisor.\\
In unsupervised
learning, there is no such supervisor and we only have input data.
The aim is to find the regularities in the input, there is a structure to the
input space such that certain patterns occur more often than others, and we want to see what generally happens and what does not. In statistics, this is called density estimation.
One method for density estimation is clustering where the aim is to
find clusters or groupings of input.\\ 
Some popular examples of unsupervised learning algorithms are:


\begin{itemize}
  \item k-means for clustering problems.
  \item Apriori algorithm for association rule learning problems.
\end{itemize}


\section{ K-Means clustering }\label{kmeans}
K-means \cite{kmeans} is one of the simplest unsupervised learning algorithms that solve the well known clustering problem. \\ K-Means clustering intends to partition n objects into k clusters in which each object belongs to the cluster with the nearest mean. This method produces exactly k different clusters of greatest possible distinction. The best number of clusters k leading to the greatest separation (distance) is not known as a priori,  and must be computed from the data. The objective of K-Means clustering is to minimize total intra-cluster variance, or, the squared error function:
\begin{figure}[H]
\centering
\includegraphics[width=0.7\textwidth]{img/ckmeans.png}
\end{figure}

 The algorithm as described by \cite{kmeans} starts with a random set of $k$ center-points ($\mu$). During each update step, all observations $x$ are assigned to their nearest center-point (see equation \ref{eqn:kmeans_assign_step}). In the standard algorithm, only one assignment to one center is possible. If multiple centers have the same distance to the observation, a random one would be chosen.

\begin{equation}
S_i^{(t)} = \big \{ x_p : \big \| x_p - \mu^{(t)}_i \big \|^2 \le \big \| x_p - \mu^{(t)}_j \big \|^2 \ \forall j, 1 \le j \le k \big\}
\label{eqn:kmeans_assign_step}
\end{equation}

Afterwards, the center-points are repositioned by calculating the mean of the assigned observations to the respective center-points (see \ref{eqn:kmeans_update_step}).

\begin{equation}
\mu^{(t+1)}_i = \frac{1}{|S^{(t)}_i|} \sum_{x_j \in S^{(t)}_i} x_j
\label{eqn:kmeans_update_step}
\end{equation}

The update process reoccurs until all observations remain at the assigned center-points and therefore the center-points would not be updated anymore.

\begin{algorithm}
\caption{Algorithm of cluster centroids initialization}

\begin{algorithmic}[1]

\Function{}{}

\State {\textbf{Input:} Data set $X = \{x^{(1)}, x^{(2)}, ..., x^{(m)} | c^{(i)} \in \mathbf{R}^n\}$}
\State {\textbf{Output:} Initial cluster centroids $\mu_{i=1, ..., k} \in \mathbf{R}^n$}
\State {Choose two samples $x^{(i)}$, $x^{(j})$ with largest distances as first two cluster centroids 
\\ \ \ \ \ \ $\mu_1 , \mu_2$ , $X \ =\ X\ -\ \{x^{(i)}, x^{(j)}\}$};

\For{$i=3, ... , k$ do }
\Comment{Choose other cluster centroids}
\State {Compute the sum of distance with all existing cluster centroids};
\State {$D^{(j)} = \sum_{t=1}^{i-1}{{\left \|x^{(j)}-\mu_t \right \|}^2}$};
\State {$X = X - x^{(j)}$};
\EndFor

\Return{$\mu_1, \mu_2, ..., \mu_k$ ;}
\EndFunction
\end{algorithmic}
\end{algorithm}


\section{Model selection } \label{ms}
Model selection is the process of choosing between different machine learning approaches,  or choosing between different hyperparameters or sets of features for the same machine learning approach  like in SVM kernels deciding between the RBF degrees and linear kernels.
The choice of the actual machine learning algorithm is less important than we'd think  there may be a "best" algorithm for a particular problem, but often its performance is not much better than other well performing approaches for that problem.


There may be certain qualities we look for in an model:

\begin{itemize}
\item Interpretable   can we see or understand why the model is making the decisions it makes?
\item Simple   easy to explain and understand
\item Accurate
\item Fast (to train and test)
\item Scalable (it can be applied to a large dataset)

\end{itemize}
Though there are generally trade offs among these qualities.\\in order to select a model from other models, we usually use one of the following approaches.\\\textbf{A Typical approach } is to take your data and split it randomly into a training set and a test set (e.g. a 70\%/30\% split). Then you train your model on the training set and see how it performs on the test set.\\The problem with this approach is that it results in an overly optimistic estimation of generalization if we tune our model's parameters with it,  so what we want to do instead is splitting the data is to not split it only into training and testing sets, but to also include a validation set. A typical ratio is 60\% training, 20\% validation, 20\% testing.

This way we can measure validation error instead of just measuring the test error .
We can use these errors   to identify what kind of problem we have if our model isn't performing well:

\begin{itemize}
\item If our training error is large and our validation/test set error is large, then we have a high bias (underfitting) problem.
\item If our training error is small and our validation/test set error is large, then we have a high variance (overfitting) problem.
\end{itemize}
As the figure \ref{fig:bias} explains : 

\begin{figure}[H]
\centering
\includegraphics[width=0.5\textwidth]{img/model.png}
\caption{ bias / variance Tradeoff}
\label{fig:bias}
\end{figure}

Some ways of evaluating a model's performance on (some of)  known data are for \textbf{validation set } to pick  the best model you can possibly get :

\begin{itemize}
\item Hold out (just set aside some portion of the data for validation; this is less reliable if the amount of data is small such that the held out portion is very small) 
\item K-fold cross-validation (better than hold out for small datasets) for better visualization  check figure \ref{fig:cross}
\begin{itemize}
\item The training set is divided into k folds
\item Iteratively take k-1 folds for training and validate on the remaining fold
\item Average the results
\item There is also "leave-one-out" cross-validation which is k-fold cross-validation where k=n (n is the number of data points)
\end{itemize}

\begin{figure}[H]
\centering
\includegraphics[width=0.9\textwidth]{img/cross.png}
\caption{K-cross validation }
\label{fig:cross}
\end{figure}


\item Bootstrapping : 
\begin{itemize}
\item New datasets are generated by sampling with replacement (uniformly at random) from the original dataset
\item Then train on the bootstrapped dataset and validate on the unselected data
\end{itemize}
\item Jackknife resampling : 
essentially to leave-one-out cross-validation, since leave-one-out is basically sampling without replacement
\end{itemize}


\subsection{Evaluating classification models}\label{eva}
Every model comes with parameters that can modify the model's behavior and performance, evaluation of the model  often comes in the form of a confusion matrix \ref{confu}.
% Please add the following required packages to your document preamble:

% Please add the following required packages to your document preamble:
% \usepackage{multirow}
\begin{table}[]
\centering
\caption{Confusion matrix}
\label{confu}
\begin{tabular}{lllclll}
\cline{3-4}
                                                                                                        & \multicolumn{1}{l|}{}           & \multicolumn{2}{l|}{\textbf{Predicted Class}}                                                                                                                                                            &                               &  &  \\ \cline{3-4}
                                                                                                        & \multicolumn{1}{l|}{}           & \multicolumn{1}{l|}{\textbf{P}}                                                                    & \multicolumn{1}{l|}{\textbf{N}}                                                                     &                               &  &  \\ \cline{1-4}
\multicolumn{1}{|l|}{\multirow{2}{*}{\textbf{\begin{tabular}[c]{@{}l@{}}Actual \\ Class\end{tabular}}}} & \multicolumn{1}{l|}{\textbf{P}} & \multicolumn{1}{l|}{\textbf{\begin{tabular}[c]{@{}l@{}}True Positives\\        (TP)\end{tabular}}} & \multicolumn{1}{l|}{\textbf{\begin{tabular}[c]{@{}l@{}}False Negatives\\        (FN)\end{tabular}}} & \multicolumn{1}{c}{\textbf{}} &  &  \\ \cline{2-4}
\multicolumn{1}{|l|}{}                                                                                  & \multicolumn{1}{l|}{\textbf{N}} & \multicolumn{1}{c|}{\textbf{\begin{tabular}[c]{@{}c@{}}False Positives\\ (FP)\end{tabular}}}       & \multicolumn{1}{c|}{\textbf{\begin{tabular}[c]{@{}c@{}}True Negatives\\ (TN)\end{tabular}}}         & \textbf{}                     &  &  \\ \cline{1-4}
                                                                                                        &                                 & \multicolumn{1}{c}{\textbf{}}                                                                      & \textbf{}                                                                                           & \multicolumn{1}{c}{\textbf{}} &  &  \\
                                                                                                        &                                 & \multicolumn{1}{c}{\textbf{}}                                                                      & \textbf{}                                                                                           & \multicolumn{1}{c}{\textbf{}} &  & 
\end{tabular}
\end{table}


The core values are:
\begin{itemize}
\item True positives (TP): samples classified as positive which were labeled positive
\item True negatives (TN): samples classified as negative which were labeled negative
\item False positives (FP): samples classified as positive which were labeled negative
\item False negatives (FN): samples classified as negative which were labeled positive
\end{itemize}


here is some important quantities we need to know in order to evaluate a model :
\begin{itemize}
\item Recall  also known as True Positive Rate (TPR): $ \frac{TP}{TP+FN}$ 
\item False Positive Rate (FPR): $\frac{TN}{TN+FP}$
\item Positive predictive value:  $\frac{TP}{TP+FP}$ 
\item Negative predictive value: $\frac{TN}{TN+FN}$
\item Precision: How many of the predicted positive samples are correctly predicted   $\frac{\text{TP}}{\text{TP} + \text{FP}}$

\item Accuracy : $\frac{TP+TN}{TP+FP+TN+FN}$

\item The F-score is often introduced as harmonic mean of precision and recall.

$F$-$score= \frac{\text{2} \cdot \text{Precision } \cdot \text{ Recall} }{ \text{Precision} + \text{ Recall}}$



\end{itemize}

\subsubsection{Area under the ROC curve (ROC AUC)}


Receiver Operating Characteristics  ( ROC) is a method that visualizes, organizes and selects models based on  their performance,  this method is for binary classification, but we can adapt it to multi-classification problem  using "OneVsAll" approach, it basically  visualizes the performance of our classifier through all the thresholds, whereas accuracy and F-score  metrics, etc, judge the performance based on particular threshold (usually  0.5) above this cut-off one label is assigned, and below it the other label is assigned. But  looking at all thresholds at once, can give a clear and more honest image of the real performance of our classifier,  since  some models may work best with a different threshold and data sometimes can be biased,   that's when  a dataset has way more or way less of the positive class than there are of the negative class. This imbalance in data can give deceptive and false results if we used accuracy as a metric, to clear the idea more let's say we have a dataset of 100 training example in which it  has 99 example positive and 1 negative example.  The accuracy here can be very deceptive with a value of 99\% so accuracy can't tell much here, the F$\textbf{-}$score can take  into consideration these kind of unbalanced data,  by considering  the true \textbf{positive } rate and the true \textbf{predictive} value  but  it does  not consider all the thresholds of the classifier. However  ROC  is insensitive to bias data ( unbalanced data ) and can run over all thresholds and plots, the  true vs false positive rates, varying the threshold can give us pair of (FPR,TPR) \ref{fig:bias2}.


\begin{figure}[H]
\centering
\includegraphics[width=0.5\textwidth]{img/roc.jpg}
\caption{ ROC curve }
\label{fig:bias2}
\end{figure}

The area under the curve (AUC) is used to quantify how good the classifier  is .  The AUC is in fact  the probability that a classifier will rank a randomly chosen positive instance higher than a randomly chosen negative one, (assuming 'positive' ranks higher than 'negative'). Generally, an AUC of above 0.8 is considered "good", an AUC of 0.5 (a straight line) is equivalent to random guessing.
And to generalize this to multi-class, assuming we have a One-vs-All (OvA) classifier, we can either go with the “micro” average or the “macro” average. In “micro averaging,” we’d calculate the performance, e.g., precision and recall, from the individual true positives, true negatives, false positives, and false negatives of the the k-class model:

$$\text{Micro-average of precision} = \frac{TP_{1} \bm{+} TP_{2} \bm{+}....\bm{+}TP_{k}}{TP_{1}\bm{+}TP_{2}\bm{+}....\bm{+}TP_{k}\bm{+}FP_{1}\bm{+}FP_{2}\bm{+}....\bm{+}FP_{k}} 
$$


$$
\text{Micro-average of recall} = \frac{TP_{1}\bm{+}TP_{2}\bm{+}....\bm{+}TP_{k}}{TP_{1}\bm{+}TP_{2}\bm{+}....\bm{+}TP_{k}\bm{+}FN_{1}\bm{+}FN_{2}\bm{+}....\bm{+}FN_{k}} 
$$


And in macro-averaging, we average the performances of each individual class:

$$\text{Macro-average precision} = \frac{PRE_{1}\bm{+}....\bm{+}PRE_{k}}{k} $$
$$\text{Macro-average recall} = \frac{RE_{1}\bm{+}....\bm{+}RE_{k}}{k} $$

Macro-average method can be used when you want to know how the system performs overall across the sets of data. 
On the other hand, micro-average can be a useful measure when your dataset varies in size.

\section{Conclusion:}

In this chapter we have presented in detail some of the most popular classification algorithms in supervised  learning,  such as SVM ( Support Vector Machine ), k-NN (k-Nearest Neighbor ) and in  unsupervised learning the  k-means algorithm, which are going to be used in the classification process after we fed them with the local descriptor, such as SURF / SIFT and feature shape descriptor as of Fourier coefficients, next we studied different methods for selecting  a model using K-fold cross validation, and their astonishing role in maximizing the performance of a classifier by choosing the optimal hyperparameters of a model [\ref{ms}].
Finally we reviewed  how to  assess the performance of a fully trained classifier using area under ROC curve ( ROC AUC ).

In the next chapter we will finally take a look on the techniques used in our HGR system, and we will discuss multiple  project results.    




\chapter{Hand gesture Recognition System} \label{hgr}

\section{introduction}

in This chapters is organized as follows, First is the Preprocessing step we will build a system that can detect and track the Hand based on minimum depth detected by Kinect sensor,then filtering any minimum depth that is not the hand, after the preprocessing  a surf and sift descriptor are applied to get description of Gestures features  and fed them into an SVM kernel via a bag of features  Approach  (also called Bag of visual Words) to reduce dimensionality of surf and sift Features space by building a codebook for our gestures, another Approach was used which is Fourier shape Descriptor where we made a dictionary for feature extraction and then Classify gestures using Nearest neighbor classifier.

\section{Data Acquisition}

to build our HGR we must first acquire Gesture's images  using a sensor to represent data from the actual world space  to the numerical space, in our system we used Microsoft Kinect Sensor\ref{sec:kinect}  to acquire both RGB and depth images fpr the  16 different gestures performed by the signer, and building a datasets of 8000 images 500  for each gesture and divided later into three partitions as we will discuss later. 



\begin{figure}[H]
\centering
\includegraphics[width=0.8\textwidth]{img/dist.png}
\caption{Depth ranges}
\label{fig:dist}
\end{figure}



as we can see in \ref{fig:dist} the depth sensor  
can give distace data until 4 meters long, but For a better resolution  the user must be  in The \textbf{optimal distance } for our HGR system Which is practically between \textbf{[80 cm, 2,5m]} of course the closer the better but it is highly recommended  to not  be any lesser than 80 cm, However we chose maximal resolution 640 $\times$ 480 to have clear details of pixels,thereafter we make a crop to the image of size (271 $\times$ 211) because it gives us enough definition to distinguish the contour  of the hands.This choice greatly reduce the number of operations, and consequently, it will improve the efficiency of the code. 

\subsection{Depth space vs color Space }
A normal color camera can be used for sign language recognition but the problem comes when we want to use it for real time applications, in this case we have to track the hand position first and then we can go for recognition but implementing the hand tracking will involve complex algorithms which will make the overall system  computationally heavy, but with the kinect depth camera this task becomes easy as it has a skeleton tracking capability by using color and depth images so we can track the approximate hand position using the skeleton data, there are 20 joints that can be tracked these joints numbers are shown in \ref{table:t1}

\begin{table}[H]
\centering
\caption{JointType Enumeration}
 \label{table:t1} 
\begin{tabu} to 0.5\textwidth { | X[l] | X[r] | }

 \hline
 SpineMid = 1  & Neck = 2    \\
 \hline 
 Head = 3  &  ShoulderLeft = 4    \\
\hline
  ElbowLeft = 5   &  WristLeft = 6    \\
\hline

 HandLeft =	7  &   ShoulderRight = 8    \\
\hline
 
ElbowRight = 9  &  WristRight = 10   \\
\hline
 
 HandRight = 11  &  HipLeft = 12   \\
\hline
 
 KneeLeft = 13  &  AnkleLeft = 14   \\
\hline
 
FootLeft = 15  &  HipRight = 16   \\
\hline
 
 KneeRight = 17  & AnkleRight = 18     \\
\hline

FootRight = 19   & SpineShoulder = 20     \\

\hline

 \end{tabu}
 \end{table}

in the following figure \ref{fig:cam9} we will show the hand position both in color and depth space using Skeletal tracking  :

\begin{figure}[H]
\centering
\includegraphics[width= 1.0\textwidth]{img/colorvsdepth.PNG}

\caption {Hand position's in both color image and depth image.
\label{fig:cam9}}
\end{figure}


By Using depth data   we eliminate all
sorts of confusion related to background. This also filters the cluttered background or overlapped images (i.e. hand overlapped with face) and illumination changes,  which justifies the use of depth sensor over RGB camera. In addition  the images of depth does not show enough contrast variations  in the hand region  so that we can get more discriminative keypoints to differentiate close gestures. 

\subsection{Vocabulary}
we performed 16 different gestures.

\begin{figure}[H]
\centering
\includegraphics[width=1.0\textwidth]{img/vocab.pdf}
\caption{ Vocabulary dataset  }
\label{fig:vocab}
\end{figure}


\textcolor{blue}{ \textbf{Note :}} our Depth datasets is public for anyone to use  the link is at the bottom of the page.\footnote{\href{https://drive.google.com/file/d/0B89YnIYeclpadHF5ajVQVUJTTjg/view}{HGR-dataset}}



\section{Hand tracking And  detection}

\subsection{Relative depth }
The First step we need to accomplish in order to start to work with the depth data, is to decide which pixels are we going to take into account to carry on the tracking. The Kinect can catch the distance of the points which are visible to the camera, between the values \textbf{minDistance and maxDistance} (\ref{fig:cam8}).

\begin{figure}[H]
\centering
\includegraphics[width=1.0\textwidth]{img/mindistance.png}
\caption{Relative depth thresholding }
\label{fig:cam8}
\end{figure}

We are going to base our choice on the object’s proximity to the Kinect. We can choose if a pixel is near based on its depth  The lesser the depth the more likely it’s the hand region.\\\textbf{ Absolute depth}: A pixel is near if its depth is lesser than a constant value, that means that the pixel is between \textbf{minDistance }and a predefined value which is greater than \textbf{minDistance }and lesser than \textbf{maxDistance}.\\\textbf{Relative depth }:First of all,  we calculate the minimum depth, and  base on that depth we select a maximum depth (by adding a constant value to the minimum depth). So, if the depth is between these two values, the pixel is near. This method allow us to have greater mobility, because it does not force us to stay in the same position the whole time. 

\textbf{In the figure \ref{fig:cam8}} the minimum depth is \textbf{minDistance} and we add a tolerance of\textbf{  t= 2cm }to it, so every point between \textbf{minDistance and minDistance + t }  will be accepted as near, that's because the hand is not flat and uniform so this tolerance will handle that problem for us.


\subsection{Hand isolation using object connectivity }

After depth thresholding we notice that there are regions that do not belong to the hand, such as what appears to be one of my  arms or some furniture or my body when my hand is closer to my body like shown in the figure  \ref{fig:cam10}. These objects just happen to be on the same depth layer as my arm and hand this means that the Distance between these objects and the Kinect lies  between the range : $$[minDistance,minDistance +2cm] \pm {Error\ committed\ by\ Kinect\ } $$.


\begin{figure}[H]
\centering
\includegraphics[width=0.5\textwidth]{img/depththresholding.png}
\caption{the depth of the body is in  the range [minDistance,minDistance +2cm]}
\label{fig:cam10}
\end{figure}

the approach used in this project  is to realize that most of the times, hands are not connected to any other object after the thresholding. We already know that the center region belongs to the hand from skeleton tracking, so we can simply apply`` floodfill ()`` offered by OpenCV to find all the connected image regions.

\textbf{ Definition and explanation of the algorithm used in floodfill () }

\begin{algorithm}[H]

\caption{Algorithm of Flood Fill }
\textbf{function}  Flood-fill(node,target\_color,replacement\_color)\\ 
\If{target\_color  ==  replacement\_color }
    {
     return
    }
\If{color of node \neq target\_color }
    {
     return
    }
Set the color of node to replacement$\_$color 

\textbf{Flood-fill}(one step to the west of node, target\_color, replacement\_color)\\
\textbf{Flood-fill}(one step to the east of node, target\_color, replacement\_color)\\
\textbf{Flood-fill}(one step to the north of node, target\_color, replacement\_color)\\
\textbf{Flood-fill}(one step to the south of node, target\_color, replacement\_color)\\
\textbf{return.} 
\end{algorithm}
with this simple algorithm we can separate objects from each other and therefor the following  Results : 
 
\begin{figure}[H]
\centering
\includegraphics[width=0.5\textwidth]{img/finalresult.png}
\caption{result after applying connectivity check for one hand  }
\label{fig:cam11}
\end{figure}

the same process done with the other hand, Making an \textbf{OR } operation to the two images we end up having two hands isolated like shown in fig \ref{fig:twohd}

 \begin{figure}[H]
\centering
\includegraphics[width=0.5\textwidth]{img/twohands.jpg}
\caption{result after applying connectivity for both hands  }
\label{fig:twohd}
\end{figure}
now after finishing  the preprocessing phase, Since we are taking a resolution of 271 $\times$ 211 Pixels, to make sure the hands are being tracked we make a bounding box that is basically related to the left hand Position, this was made under the assumption that sign language  usually performed with two hands close to each other, \\
bounding box steps : 
\begin{itemize}
\item  Use LEFT hand location provided by SDK Skeleton  tracker (Px, Py)
 \item Crop to box around hand using depth at (Px, Py)
    \item scale = $\frac{20000.0f}{Depth(px,py)}$ 
    \begin{enumerate}
    \item Constant 20000.0f  found empirically
    \item Proportional to the area of bounding box 
    \end{enumerate}
\item Top left Pixel  = ( Px -10 * Scale * 0.5, Py - 10*Scale * 0.5)
\begin{enumerate}
    \item This formula scales an initial 10x10 box centered at (Px, Py) to
the appropriate dimensions based on the depth 
    \item Constants 0.5  take into account that most hand
silhouettes are more tall than wide 
    \end{enumerate}
   \item Width = Height = 10 * Scale 
\end{itemize}

\textcolor{blue}{\textbf{Summary of the section in below picture :} }

 \begin{figure}[H]
\centering
\includegraphics[width=18cm, height=16cm ]{img/preprocessing.pdf}
\caption{result after applying connectivity for both hands  }
\label{fig:preprocessing}
\end{figure}

\textbf{ motion capturing of hand articulation is a challenging task, since the hand presents a motion of high degrees of freedom (rotation variance). Additionally, Our system is based on relative depth which allow user to move around the camera This throws a problem of scale variance}, To solve these issues We chose the best feature descriptor that are most suitable for such a task,which are  SIFT AND SURF descriptor that are both Scale invariant and rotation invariant.

in the next section we will build the recognition System using these descriptors.

\section{ Gesture recognition using surf and sift}


\subsection{Feature Descriptor selection :}
in order to Choose the most suitable descriptor for our HGR system we must choose wisely what kind of descriptor can give us the most promising results, as such  HOG(histograms of oriented gradients) feature that  has been introduced into \textbf{gesture recognition }, but the HOG features is not ideal for gesture recognition for it's computation in dense grids at some single scale without orientation alignment. SIFT ( Scale Invariant Feature Transform) has shown good achievement over the years in object recognition and thus is a good candidate for our system.\\ The contribution of this project is that we propose surf and sift feature descriptors of hand gesture images and then bag of visual words are to map these descriptors to a dimension vector and a  support vector machine(SVM) classifier is then  trained to recognize hand gestures.
\begin{enumerate}
\item We show that this model is not only effective for common hand gestures dataset, but also achieves good  recognition rate for depth data collected by Kinect sensor.
\item Hand gestures training images can be represented by sets of keypoint descriptors generated by SIFT or SURF, but the numbers of keypoints from the images are \textbf{different and lack meaningful ordering }. This creates difficulties for machine learning methods. To address  this problem, we use the bag-of-words ( also called bag of features ) approach.

\end{enumerate}

\subsection{Bag Of visual Words Model}

\begin{wrapfigure}{r}{0.28\textwidth}
\includegraphics[width=0.25\textwidth]{img/codebook.jpg}
\caption{Example of Visual words in a codebook}
\label{fig:coodbook}
\end{wrapfigure}

\textbf{Bag-of-words} (also called bag of features) is one of the most famous approaches in computer vision, the idea is almost like a  dictionary of local interest points (keypoint ) of an image or images’ datasets. where the object is  represented as a bag of visual words fig [\ref{fig:coodbook}].
Basically, the important points in images are the visual words. These points are features  and they are discriminative and very stable to rotations changes and scale changes. 
The BoW model can be used for image classification, by building a vocabulary database(codebook) of many visual words and represent every image as a histogram of the frequency words that  are in  the image. In order to use BoW for image classification, these features need to be extracted and then we need to generate a codebook, and a histogram afterwards. To retrieve these features,we use some of the local features descriptors such as (SIFT, SURF,HOG,GLOH,ORB,FAST..).
Following this step, each image is a collection of  vectors of the same dimension, \textbf{in which the order of different vectors is not important, but the sets diversify in cardinality and lack significant ordering}. This creates difficulties for learning methods (e.g., classifiers) that  require feature vectors of  \textbf{fixed dimension}  as input.
the vector quantization technique regroups ( clusters)  the keypoint descriptors in their feature space into large sets using the clustering algorithm K-means and encodes each keypoint to the index of the cluster to which it belongs. Each cluster is a visual word that represents a specific local pattern shared by the keypoints in that cluster. Thereby the clustering process generates vocabulary of visual words describing local patterns in images. The codebook size is determined by the number of clusters, therefore, each patch in an image is mapped to a particular codeword.


\subsection{keypoint descriptor Clustering }

\begin{enumerate}
\item \textbf{Feature extraction and description.} It is to extract a representative local features from an image as a description (surf or sift ).
\item \textbf{ Dictionary}
using K-means to cluster Surf or Sift features into K centroids this process is detailed in section [\ref{kmeans}],  Once Clustering is done, each feature vector (keypoint) is assigned to one and only one cluster center that is in the nearest distance with respect to the Euclidean metric in 64 / 128 dimensional feature vectors. The keypoints that are assigned to the same cluster center will be in the same subgroup so that after clustering we have k disjoint subgroups of keypoints. Therefore k-means clustering decreases dimensionality for every training image with n keypoints (n $\times$ m)) to 1 $\times$ k, where m is number of feature descriptor can be either 63 or 128, k is number of clusters.
\item \textbf{ Generating bag-of-features} in this process we generate a bag of features for input images, That means, each keypoint, extracted from input  image, will be represented by one component in the generated bag-of-feature vector with value equal to the index of the
centroid in the cluster model with nearest Euclidean distance. Figure \ref{fig:histo} shows the process of generating the bag-of-features 

\end{enumerate}


\subsection{classification algorithm }

\textbf{Training Processes are as follows }:


\begin{enumerate}
 \item extracting SURF and SIFT features from training samples with different scales and  orientations.
 \item applying  K-means clustering on the training images will help build the Visual words each visual words is considered as centroid of a subgroup of the same Features, Now we can generate bag of features ( histogram of visual words ) for training images where each keypoint will be extracted from training image, will be represented by one component in the generated bag-of-feature vector with value equal to the index of the centroid in the cluster model with nearest Euclidean
distance this basically means getting frequencies of Visual words in the Current image, to have better visualization  about this step see the Figure \ref{fig:histo}

\begin{figure}[H]
\centering
\includegraphics[width=18.5cm, height=11cm ]{img/bow/transparantnewhistogram.pdf}
\caption{ generating histogram for training image and test image by Mapping keypoint into Bag-of-Features space }
\label{fig:histo}
\end{figure}

\item SVM uses a non-linear mapping to transform original training dataset into higher dimension and searches for a linear optimal
separating hyperplane. since our Gesture can have the same shape and certainly often have the same features for two different gestures  this will make Linear separation Hard and can results in false results the kernel method can be used in our case  to transform data  in to higher dimension space in which they are separable. This process will allow classification of multiple classes (16 class ) 
\end{enumerate}


\textbf{validation :}
we used K-fold cross validation method to tune the parameters of our SVM, Using both linear and RBF kernels, Also K that represents Number of clusters plays a role in tuning our model, the result will be shown and discussed in the next section \ref{resultssvm} 



\textbf{Classification process: }
\begin{enumerate}
    \item Extract Keypoint using surf and sift to the test images.
    \item  representing test images by Histogram
    \item using the trained SVM classifier to obtain a classification result.
    
\end{enumerate}

the following Figure \ref{fig:algo1}  summarize the whole  process :

\begin{figure}[H]
\centering
\includegraphics[width=17cm, height=17cm ]{img/bow/mynewalgo.pdf}
\caption{Flowchart for the structure of the algorithm, in block (A) the Extracted Keypoints  are Clustered to i Vocabulary and generate a codewords dictionary,then compute a histogram using visual words ( vocabulary ) for all the training samples these histograms are then fed to an SVM for classification, in Block (B) test images are converted to histogram representation by counting how many keypoints appeared in cluster i,  and then classify the results using the trained SVM }
\label{fig:algo1}
\end{figure}




\subsection{SIFT and SURF  Experimental Results}\label{resultssvm}
in this section we are interested in maximizing the performances of SVM, using K-cross Validation to make model selection, with  K=10 for cross Validation which has ensure us great tuning for our models parameters,
we have chose AUC of ROC curve  to be our main Evaluation Metric since The method can be used for binary classification, therefor we used \textit{one versus all} method to convert the problem to a binary task.

\textbf{Dataset division :}

we took 500 image for each gesture, we divided our dataset into three partitions, training set, validation set,and Test set.

\begin{itemize}
    \item training set :
         300 images for each gestures. with total of 4800 image
    \item validation set :
         100 images for each gestures. with total of 1600 image
    \item test set :
         100 images for each gestures.with total of 1600 image
\end{itemize}



\subsubsection{ SVM performance variation in function of the codebook's size }
The codebook size will be determined depending
on the type  of the data processed. There will be a sort of
compromise  for choosing number of clusters (size of codebook). If it is too large, then there will be quantization artifacts and overfitting because of insufficient samples of the keypoints extracted from the training image. If it is too small, then each bag of words vector will not represent all the keypoints extracted from its related image.

since the structure of our data is unique and has not a lot of keypoint variations, as a matter of fact the hands are the only object detected in all the scenes. so we will start with a  small number for k  and change it until finding a reasonable  AUC ( note : size of codebook is in fact the number of clusters used in K-means).

\begin{figure}[H]
\centering
\includegraphics[width=18cm, height=5cm]{img/bow/dictionnary1.png}
\caption{}
\label{fig:codbooksize}
\end{figure}

\subsubsection{SVM linear kernel}

the feature vector resulted  by surf and sift  in the training set is  519523 for SURF  and 208397 for SIFT, using  bag of feature we reduce it to (4800 $\times$ 500) where 4800 is the number of all training set and 500 is the number of codewords in our codebook  that represent for SVM the number of features, even with dimensionality reduction it is still a big feature vectors to train, thus the choice of linear kernel that generate a limited number of support vectors.


\subsubsection{Surf performance: }
first we varied the Hyper-parameter \texttt{C}, in order to to find an optimal value that maximize SVM performance, the following figure \ref{fig:C} illustrate the variation  of the average AUC of ROC in function of C variations.

\begin{figure}[H]
\centering
\includegraphics[width=0.56\textwidth]{img/C_values.png}
\caption{the choice of hyper-parameter C }
\label{fig:C}
\end{figure}

C $=$ 100, proved to be the best for linear kernel that scored 90\% Average AUC, thus it was chosen as our Hyperparameter for our model.


the next figure is a representation of ROC curves of Linear Kernel with different C :

\begin{figure}[H]
\centering
\includegraphics[width=16.5cm, height=11cm]{img/roc/LINEAR_svm_onevsal_ROC_class_C100.png}
\caption{ROC of 16 gestures with  C = 100}
\label{fig:roc}
\end{figure}

\textbf{interpretation }
as shown in the figure above the average AUC is equal to 90 \%
that represent the  overall estimation of our classifier to rank positive data ( in our case is the gesture X $\in$ [0,15]) higher than a negative data ( the rest of gestures excluding the chosen gesture X ),  for instance the gesture 11 that corresponds to gesture V  that has an AUC equal to 99\% this means the classifier can perfectly classify V  among the rest gestures, however the classifier can perform poorly in the case of gesture 14 ("M")with an AUC of  50\%  this means the classifier can  expect  as many gesture M  as other gestures. 

\subsubsection{Sift performance : }
As for sift the optimal value for C is  1.0, to find the best hyperparameter that maximize the linear kernel we varied C and calculates the average ROC AUC fig [\ref{fig:Csift}] :

\begin{figure}[H]
\centering
\includegraphics[width=0.6\textwidth]{img/roc/sift_C_values.png}
\caption{the choice of hyper-parameter C }
\label{fig:Csift}
\end{figure}

Next we present ROC curve of 16 gestures with  C= 1.0 using Sift features. 

\begin{figure}[H]
\centering
\includegraphics[width=16cm, height=12cm]{img/roc/sift_ROC_class_C.png}
\caption{ROC of 16 gestures with  C = 1.0 }
\label{fig:rocsift}
\end{figure}

\textbf{interpretation :}
The average AUC is 88\%, under this performance the class  13,14,4 that corresponds to gestures ( "M", "O", "E" ), this means the classifier can give random guessing and we expect it to classify "M","O","E" as equally as to all the other gestures. But the classifier was able to achieve good result in almost  all other gestures.


\subsubsection{SVM RBF kernel}

for RBF kernel, we have two parameters to tune our model with, the hyperparameter C and the parameter $\gamma$. we have varied the C in the interval of [$10^{-3}$,0.01,...,$10^{7}$] and $\gamma$ in [$10^{-5},....,10^{3}$], the best parameters were found are C $=$ 1000.0, $\gamma$ $=$ [0.1,1.0,0,10]. these values gave a performance of 98\%.

\subsubsection{Surf performance : }

\begin{figure}[H]
\centering
\includegraphics[width=0.71\textwidth]{img/grid/grid_searchbone.png}
\caption{Gridsearch for optimal C and $\gamma$}
\label{fig:grid}
\end{figure}



the following figure is a representation of ROC curve for C $=$ 1000.0, and $\gamma$ $=$ 0.1.

\begin{figure}[H]
\centering
\includegraphics[width=16.5cm, height=10cm]{img/roc/RBF_C1000_gamma01.png}
\caption{Roc for RBF kernel with C= 1000.0 and gamma = 0.1}
\label{fig:rbf0}
\end{figure}
the average AUC is 98\% this describe the probability to classify a gesture X $\in$ \{ gestures\} against other gestures $\in$ \{\{gestures\} - \{X\} \} is almost ideal.
Now let's see the  ROC curve for C= 1000.0 and $\gamma$ = 1.0, that gave an Average AUC of 98 \% as well.
\begin{figure}[H]
\centering
\includegraphics[width=16.5cm, height=11cm]{img/roc/RBF_C1000_gamma10.png}
\caption{Roc for RBF kernel with C= 1000.0 and gamma = 1.0}
\label{fig:rbf1}
\end{figure}

in fact the best optimal gamma was between [0.1,1.0], we have choose C$=$ 1000.0, $\gamma$ = 1.0 for our RBF model.

\subsubsection{Sift perforamnce :}



the following figure is a representation of ROC curve for C $=$ 1000.0, and $\gamma$ $=$ 0.1.

\begin{figure}[H]
\centering
\includegraphics[width=16.5cm, height=11cm]{img/roc/sift_ROC_class_C1000gamma.png}
\caption{Roc for RBF kernel with C= 1000.0 and gamma = 0.1}
\label{fig:rbf2}
\end{figure}


\subsubsection{comparing results  } 

 We found that SURF feature descriptors are more appropriately for visual
hand gesture recognition than SIFT feature descriptors.
SIFT is  slow which  gave lag during feature extraction. on the other hand SURF performs well on Depth and blurry images and is typically 2 to 3 times faster than SIFT as shown in the table below. We took 4800 images of 200 images from each class ( 16 $\times$ 300 ) and registered  the time required for extraction and description  by SURF and SIFT.
\begin{table}[!h]
\centering

\begin{tabular}{c|c|c|ll}
\cline{2-3}
\multicolumn{1}{l|}{\textbf{}}                                                                                                                      & \textbf{SIFT}        & \textbf{SURF}                            &  &  \\ \cline{1-3}
\multicolumn{1}{|c|}{\textbf{\begin{tabular}[c]{@{}c@{}}Average time (in seconds ) consumed for \\ features detection and description Per Image \end{tabular}}} & \textbf{0.30s} & \multicolumn{1}{l|}{\textbf{0.12s}} &  &  \\ \cline{1-3}
\multicolumn{1}{|c|}{\textbf{number of features extracted per image }}                                                                                         & \textbf{43}        & \textbf{107}                            &  &  \\ \cline{1-3}

\end{tabular}
\caption{Average speed time consumed for features detection and extraction And accuracy for SURF and SIFT applied on a set of 4800 images for 16 different gestures }
\label{siftvssurfacc}
\end{table}

Surf Can be 3  times faster than sift also Surf Was capable of describing Gesture keypoint better than sift by difference of 50 keypoints  it can gives an Idea about the descriptor that suits best our data and can extract as much information as needed for recognition. The average ROC AUC  in sift is 88\% in linear kernel while surf scored 90\%, and as for RBF kernel still surf was the winner by 98\%  against 91\%.
For all the Reasons cited above  we found SURF descriptors to be apt for our application to meet the requirement  \textbf{of a  Real time application} since Sift descriptor can not be used in real time applications however Not with normal machines.

%%%%%%%%%%%%%%%%%%%%%%%%%%%%%%%%%%%%%%%%%%%%%%%%%%%%%%%%%%%%%%%%%%%%%%%%%%%%%%%%%%

\section{gesture recognition using Fourier Descriptor and NN  classifier }

 The use of depth image has benefit
in segmenting hand image rather than color-based
segmentation, segmentation in depth image is more robust
since the lighting variation does not affect the image quality.
 Fourier descriptors represent feature of each hand gesture as Unique signature As long as the shape (images ) are different from each others.
Generally, this section is separated into two phases:
dictionary build phase and classification phase. \\In the
dictionary build phase, the Fourier descriptors of each
character are stored into a database to develop gesture
dictionary. The gesture dictionary and comparison methods are
derived from \cite{clif}.\\ The classification phase has the similar
step, except that the Fourier descriptors are compared with the
dictionary using distance metric as classification methods.
The result of classification phase is the meaning of the
acquired gesture sign.
\texttt{For more details about the algorithm used for this project please check the section \ref{FDT}
}

\subsection{Classification }
Now we have a set of features for each image and the next
step is to train a classifier to classify the different hand shapes
used in testing. the classifier used in for Fourier descriptor is based on Euclidean distance which is a very simple way for
classification.\\
Classification using Euclidean Distance: The set of
features induces a distance on the shape space, which is given
by the Euclidean distance: 
$$d(f_{1},f_{2}) = \sqrt{\sum_{k=0}^{N-1} |f_{1}(k)-f_{2}(k)|^{2}}$$\\
\texttt{ Where $f_{1},f_{2}$ are the feature vectors for the two images being
compared. The two vectors with the least distance will have
the same class}\\

The following Figure \ref{fig:dft} summarize the whole process :  
        
\begin{figure}[H]
\centering
\includegraphics[width=0.8\textwidth]{img/latex.png}
\caption{the architecture used for FDs}
\label{fig:dft}
\end{figure}




\subsection{Fourier descriptor Experimental Results :}

For each shape we select a set of points with the equal point
sampling method. To use DFT we choose a number that is a
power-of-two. We tested the classifiers with initial 128, 64, 32, 16
and 8 points. Those features are used as a training set to be
used in classification. To classify the test image we used the Classifier Nearest Neighbor that uses simple  Euclidean distance, where we calculated the distance between the features of each image in
the test set with the features of the Dictionary.


\subsubsection{Accuracy variation in function of number of Fourier descriptor features : }


// still in the process..


\section{comparing results for FD}




...



\section{Conclusion}


........
\chapter*{Conclusion and Perspective}

In this work, we succeed in building an HGR system that is able to detect the presence of two hands, to recognize the meanings of a gesture  in a predefined vocabulary. The lighting conditions, signers skin colours and clothing, and background  have little impact on the performance of this HGR system. 
In this project we used Microsoft kinect for data acquisition of depth images. At  First we were focusing on giving the user a better experience by tracking and  detecting the hand  in a range of [80cm , 3 meters]  allowing the  user to move freely around the camera. then we extract features using two powerful descriptors SIFT(Scale Invariant Features Transform) and SURF(Speeded Up Robust Features) a  third descriptor(Fourier shape descriptor)  was added to complete the requirement of a robust HGR system. having our stable keypoints that are invariant to scale and rotation  we used machine learning algorithms SVM, K-means and Nearest Neighbors  to classify our  16 gestures. afterwards we generate a codebook based on the keypoint of sift and surf,we also did  a comparative study between sift and surf performance's based on their generalization performance, Finally we conclude that surf outperform sift with an Average AUC of 98\% with RBF kernel against 91\% for sift, the only problem with these two descriptors is that they are computationally expensive which blocks their application on rela time applications,Hence the need for a fast descriptor such as fourier descriptor.

\begin{itemize}
    \item this application can detect  in a minimum range of 80 cm and maximum range of 2 and half meters until 3 meters.
    \item The application is optimized and make no calculation or classification if there is no \textbf{Human body} in front of the camera, the program starts only when u make a gesture to recognize the human skeleton.
    \item A powerful CPU or  GPU can make  sift and surf descriptors along with BOF algorithm to work fine especially with the recent implementations of OpenCV for SURF GPU and SIFT GPU.
    \item the project is extendable in term of adding any new gestures to the vocabularies (as long as  they have a different shape).
    \item Kinect Version 1 (used for this project) gives limited depth range of 3 meters (theoretically 4 meters),and gives false hands positions  especially if the hands cross. he error increases respectively with depth.  To the author  best knowledge the Kinect version 2 gives more depth range until 8 meters and uses good filters for to stabilize body joints positions. 
    \item the hands images taken by the camera were only of an individual person, which is the author of this project, this can probably causes our application to overfit a little bit, but this problem can be easily solved with if we shoot with  the different persons with different shape of hands and to account for the different ways that people might perform those gestures.this will help a lot in boosting the generalization performance of our application.
   
\end{itemize}


This project was an opportunity to apply the different knowledge acquired during our studies in the IST master; especially as we have deepened our knowledge in the field of computer vision by applying image processing techniques using the OPENCV library. We also learned new skills at programming with the Kinect SDK. Moreover; this work allowed us to master some of Machine learning aspects as well as learning the necessary steps to have a good performance of generalization.


Our application performs the recognition of the static gestures of the both  hands. Thereby It is very important to deepen this work by recognizing the gestures dynamically in order to establish a system that is able to translate the sign language into a natural language and include Text To Speech (TTS) protocol. 





\appendix
\include{appendices/appendiceKinect}
\include{appendices/appendiceOpencv}



%%%%%%%%%%%%%%%%%%%%%%%%%%%%%%%%%%%%%%%%%%%%%%%%%%%%%%%%%%%%%%%%%%%%%%%%%%%%%%%%%%%%%%%%%%%%%%%%%%%%%%%%%
\medskip



\begin{thebibliography}{80}


\bibitem{20}D. S. Zhang and G. J. Lu. A Comparative Study on Shape Retrieval
Using Fourier Descriptors with Different Shape Signatures. In Proc. Int.
Conference on Multimedia and Distance Education. Fargo, ND, USA,
pp.1-9, June 2001. 
\bibitem{21}
Kai Guo; Ishwar, P.; Konrad, J., "Action Recognition in Video by Covariance
Matching of Silhouette Tunnels," Computer Graphics and Image Processing
(SIBGRAPI), 2009 XXII Brazilian Symposium on , vol., no., pp.299,306, 11-15 Oct.
2009 


\bibitem{27} H. Cooper, B. Holt, and R. Bowden, “Sign language recognition,” in in Visual Analysis of
Humans, Springer, 2011, pp. 539–562. 



\bibitem{29} T. Takahashi and F. Kishino, “Hand gesture coding based on experiments using a hand
gesture interface device,” ACM SIGCHI Bull., vol. 23, no. 2, pp. 67–74, 1991. 



\bibitem{31}T. K., “Glove lends the deaf a hand,” USA Today, 2002. [Online]. Available:
http://www.usatoday.com/news/science/2002-01-16-deafglove.htm. [Accessed: 01-Dec2013].

\bibitem{32} T. Starner and A. Pentland, “Real-time american sign language recognition from video
using hidden markov models,” in in Motion-Based Recognition, Springer, 1997, pp. 227–
243. 

\bibitem{33}
 T. Starner, J. Weaver, and A. Pentland, “Real-time american sign language recognition
using desk and wearable computer based video,” Pattern Anal. Mach. Intell. IEEE Trans.,
vol. 20, no. 12, pp. 1371–1375, 1998. 

\bibitem{35}
Z. Ren, J. Meng and J. Yuan, "Depth camera based hand gesture recognition and its applications in Human-Computer-Interaction," 2011 8th International Conference on Information, Communications and Signal Processing, Singapore, 2011, pp. 1-5.

\bibitem{36} P. Vijay and N. Suhas, “Recent Developments in Sign Language Recognition: A Review,”
irdindia.in, no. 2, pp. 21–26, 2012. 
 
\bibitem{39}M. Majid and J. Zain, “A REVIEW ON THE DEVELOPMENT OF INDONESIAN SIGN
LANGUAGE RECOGNITION SYSTEM,” J. Comput. Sci., vol. 9, no. 11, pp. 1496–
1505, 2013. 

\bibitem{krig}
Krig S. Interest Point Detector and Feature Descriptor Survey. In: Computer Vision Metrics. Springer; 2014. p. 217–282.



\bibitem{lowe}
Lowe, David G. (1999). "Object recognition from local scale-invariant features" (PDF). Proceedings of the International Conference on Computer Vision. 2. pp. 1150–1157. doi:10.1109/ICCV.1999.790410.



\bibitem{47} D. Capilla, “Sign Language Translator using Microsoft Kinect XBOX 360 TM,” 2012. 


\bibitem{dardas}
N. Dardas, Q. Chen, N. D. Georganas and E. M. Petriu, "Hand gesture recognition using Bag-of-features and multi-class Support Vector Machine," 2010 IEEE International Symposium on Haptic Audio Visual Environments and Games, Phoenix, AZ, 2010, pp. 1-5.


\bibitem{MM} Tuytelaars, T., Mikolajczyk, K.: Local invariant feature detectors: a survey. Found. Trends
Comput. Graph. Vis. 3(3), 177–280 (2007)




\bibitem{f} Crowley, James L., and Alice C. Parker, “A representation
for shape based on peaks and ridges in the difference of lowpass transform,”
IEEE Trans. on Pattern Analysis and Ma-
chine Intelligence,
6, 2 (1984), pp. 156–170.

\bibitem{e}Lindeberg, Tony, “Detecting salient blob-like image structures and their scales with a scale-space primal sketch:
a method for focus-of-attention,”
International Journal of
Computer Vision,
11, 3 (1993), pp. 283–318.



\bibitem{h} Oliva, A., Torralba, A.: Modeling the shape of the scene: a holistic representation of the spatial
envelope. Int. J. Comput. Vis. 42(3), 145–175 (2001)
\bibitem{i} Bianco, S., Mazzini, D., Pau, D., Schettini, R.: Local detectors and compact descriptors for
visual search: a quantitative comparison. Digital Signal Process. 44, 1–13 (2015)
\bibitem{k} Jégou, H., Perronnin, F., Douze, M., Sánchez, J., Pérez, P., Schmid, C.: Aggregating local
descriptors into a compact codes. IEEE Trans. Pattern Anal. Mach. Intell. 34(9), 1704–1716
(2012)
\bibitem{kinect}
Microsoft, “Kinect for Windows Sensor Components and Specifications.” [Online]. Available: \url{http://msdn.microsoft.com/en-us/library/jj131033.aspx}Accessed: 06-Dec-2013 ]

\bibitem{kinect1}
 Microsoft, “Kinect Fact Sheet,” 2010. [Online]. Available: \url{www.microsoft.com/enus/news/presskits/xbox/docs/kinectfs.docx}.
\bibitem{kinect15}
M. Andersen, T. Jensen, and P. Lisouski, “Kinect depth sensor evaluation for computer
vision applications,” 2012.
\bibitem{kinect17}
Microsoft, “Kinect for Windows - Human Interface Guidelines.” [Online]. Available:
\url{http://msdn.microsoft.com/en-us/library/jj663791.aspx}.
\bibitem{kinect18}
J. Webb and J. Ashley, Beginning Kinect Programming with the Microsoft Kinect SDK.
Apress, 2012
\bibitem{kinect20}
Microsoft, “Coordinate Space.” [Online]. Available: \url{http://msdn.microsoft.com/enus/library/hh973078.aspx#Depth\_Ranges}.
[Accessed: 06-Dec-2013].

\bibitem{clif}
C. Chan, S. S. Mirfakhrae, “ Hand Gesture Recognition using
Kinect”, Bchelor Thesis, Boston University Department of
Electrical and Computer Engineering ,Boston, Dec 13, 2013.
\bibitem{kmeans}
J. B. MacQueen (1967): "Some Methods for classification and Analysis of Multivariate Observations, Proceedings of 5-th Berkeley Symposium on Mathematical Statistics and Probability", Berkeley, University of California Press, 1:281-297
\bibitem{Arthur} 
Arthur, D.; Vassilvitskii, S. (2007). "k-means++: the advantages of careful seeding" (PDF). Proceedings of the eighteenth annual ACM-SIAM symposium on Discrete algorithms. Society for Industrial and Applied Mathematics Philadelphia, PA, USA. pp. 1027–1035.

\bibitem{svm}
Hill, T.,  Lewicki, P. (2006). Statistics Methods and applications : a comprehensive reference for science, industry, and data mining. Tulsa, Okla: StatSoft.

\bibitem{wfd}
 "FAQ". World Federation of the Deaf. Retrieved 17 August 2013.
 
 
\end{thebibliography}



\end{document}
